\documentclass[conference]{IEEEtran}
\usepackage{booktabs}
\usepackage{graphicx}
\usepackage{amsmath}
\usepackage{url}

\title{How Much Does Codec-Aware Augmentation Help Lightweight Audio Deepfake Detection?\\
\large An Empirical Study on Codec-Degraded and Telecom Speech}

% TODO before submission: verify venue template (IEEEtran for ICASSP/SLT/SPL;
% Interspeech uses its own single-column template).

\author{\IEEEauthorblockN{Harshil, Madhigank, and Dr. Chetali}
\IEEEauthorblockA{Department of AI \& ML, <College Name>}}

\begin{document}
\maketitle

\begin{abstract}
TODO (write last). One paragraph: problem, gap (synthetic codec augmentation
understudied with TRUE codecs; generalization of lightweight CMs to codec
degradation), method (lightweight CNN; five augmentation strategies incl.\ true
Opus encode/decode vs.\ naive resampling; 3 seeds; official metrics incl.\
min t-DCF with provided ASV scores), key findings (numbers TBD from the grid),
significance.
\end{abstract}

\begin{IEEEkeywords}
audio anti-spoofing, deepfake speech detection, data augmentation, codec
robustness, ASVspoof
\end{IEEEkeywords}

\section{Introduction}
TODO. Motivation: voice deepfakes over phone/voice-chat channels (WhatsApp,
Teams) are the realistic attack surface; classifiers trained on studio-quality
audio degrade under codecs. Contribution list (mirror paper/PLAN.md):
\begin{itemize}
\item systematic ablation of five augmentation strategies (incl.\ TRUE Opus vs.\
naive resampling) on a lightweight CNN, 3 seeds, official metrics;
\item a codec robustness matrix: models evaluated under 11 real codec conditions
(Opus/MP3/AAC/AMR-NB at multiple bitrates);
\item cross-corpus evaluation on ASVspoof 2021 LA (real telecom channels);
\item official min t-DCF with the provided ASV scores and per-attack analysis.
\end{itemize}

\section{Related Work}
TODO. Seed list in references.bib: ASVspoof 2019/2021/5, t-DCF (Kinnunen et
al.), LFCC-GMM/LSTM baselines, RawNet2, RawBoost, AASIST, wav2vec2-based CMs,
In-the-Wild generalization (M\"uller et al.), telephony spoof detection.

\section{System and Augmentation Strategies}
\subsection{Baseline countermeasure}
3-block CNN on $128\times128$ log-mel spectrograms (16 kHz, 4.065\,s crops,
hop 512), $\approx$X\,M parameters. Global-stat normalization. No class
weights. Adam, lr $2\cdot10^{-4}$, max 60 epochs, model selection on dev EER.

\subsection{Augmentation modes}
none / noise (SNR 5--20 dB) / resample (16k$\to$8k$\to$16k decimation) / opus
(true libopus at \{6,12,24,48\} kbps) / full (50\% none, 25\% opus, 25\% noise).
Static pool of $K{=}8$ degraded variants per utterance; details in the repo.

\section{Experimental Setup}
Datasets: ASVspoof 2019 LA (train 25{,}380 / dev 24{,}841 / eval 71{,}237);
ASVspoof 2021 LA eval for cross-corpus. Metrics: EER and official minimum
normalized t-DCF (ASV at its own EER operating point; provided ASV scores).
Three seeds (42/123/456), mean $\pm$ std, bootstrap CIs on EER.

\section{Results}
\input{tables/ablation}          % generated by scripts/stats.py
\input{tables/codec_matrix}      % generated by scripts/stats.py
\input{tables/per_attack}        % generated by scripts/stats.py
TODO: training curves figure; efficiency table (params, MACs, latency);
ASVspoof 2021 LA cross-corpus table; discussion of unknown-attack behavior.

\section{Conclusion}
TODO.

\bibliographystyle{IEEEtran}
\bibliography{references}
\end{document}
